% concatenated from Triangular.tex
%Triangular
%June 24, 2025-submitted to Integers 


 


\documentclass{amsart}


\usepackage{amsmath,amssymb,amsthm,latexsym}






\newtheorem{theorem}{Theorem}
\newcommand{\bt}{\begin{theorem}}
\newcommand{\et}{\end{theorem}}
\newtheorem{lemma}{Lemma}
\newcommand{\bl}{\begin{lemma}}
\newcommand{\el}{\end{lemma}}
\newtheorem{corollary}{Corollary}
\newcommand{\bc}{\begin{corollary}}
\newcommand{\ec}{\end{corollary}}
\newcommand{\bconj}{\begin{conjecture}}
\newcommand{\econj}{\end{conjecture}}
\newtheorem{problem}{Problem}
\newcommand{\bprob}{\begin{problem}}
\newcommand{\eprob}{\end{problem}}
\newcommand{\beq}{\begin{equation}}
\newcommand{\eeq}{\end{equation}}
\newcommand{\benum}{\begin{enumerate}}
\newcommand{\eenum}{\end{enumerate}}
\newcommand{\N}{\ensuremath{ \mathbf N }}
\newcommand{\Z}{\ensuremath{\mathbf Z}}
\newcommand{\Q}{\ensuremath{\mathbf Q}}
\newcommand{\R}{\ensuremath{\mathbf R}}
\newcommand{\Rn}{\ensuremath{\mathbf R^n}}
\newcommand{\C}{\ensuremath{\mathbf C}}
\newcommand{\mbA}{\ensuremath{ \mathbf A}}
\newcommand{\mbB}{\ensuremath{ \mathbf B}}
\newcommand{\mbC}{\ensuremath{ \mathbf C}}
\newcommand{\mbF}{\ensuremath{ \mathbf F}}
\newcommand{\PP}{\ensuremath{ \mathbf P}}


\newcommand{\mca}{\ensuremath{ \mathcal A}}
\newcommand{\mcb}{\ensuremath{ \mathcal B}}
\newcommand{\mcc}{\ensuremath{ \mathcal C}}
\newcommand{\mcd}{\ensuremath{ \mathcal D}}
\newcommand{\mce}{\ensuremath{ \mathcal E}}
\newcommand{\mcf}{\ensuremath{ \mathcal F}}
\newcommand{\mcg}{\ensuremath{ \mathcal G}}
\newcommand{\mch}{\ensuremath{ \mathcal H}}
\newcommand{\mci}{\ensuremath{ \mathcal I}}
\newcommand{\mcj}{\ensuremath{ \mathcal J}}
\newcommand{\mck}{\ensuremath{ \mathcal K}}
\newcommand{\mcl}{\ensuremath{ \mathcal L}}
\newcommand{\mcm}{\ensuremath{ \mathcal M}}
\newcommand{\mcn}{\ensuremath{ \mathcal N}}
\newcommand{\mco}{\ensuremath{ \mathcal O}}
\newcommand{\mcp}{\ensuremath{ \mathcal P}}
\newcommand{\mcr}{\ensuremath{ \mathcal R}}
\newcommand{\mcs}{\ensuremath{ \mathcal S}}
\newcommand{\mct}{\ensuremath{ \mathcal T}}
\newcommand{\mcu}{\ensuremath{ \mathcal U}}
\newcommand{\mcv}{\ensuremath{ \mathcal V}}
\newcommand{\mcw}{\ensuremath{ \mathcal W}}
\newcommand{\mcx}{\ensuremath{ \mathcal X}}

\newcommand{\mba}{\ensuremath{ \mathbf a}}
\newcommand{\mbb}{\ensuremath{ \mathbf b}}
\newcommand{\mbc}{\ensuremath{ \mathbf c}}
\newcommand{\mbe}{\ensuremath{ \mathbf e}}
\newcommand{\mbm}{\ensuremath{ \mathbf m}}
\newcommand{\mbo}{\ensuremath{ \mathbf 0}}
\newcommand{\mbu}{\ensuremath{ \mathbf u}}
\newcommand{\mbv}{\ensuremath{ \mathbf v}}
\newcommand{\mbw}{\ensuremath{ \mathbf w}}
\newcommand{\mbx}{\ensuremath{ \mathbf x}}
\newcommand{\mby}{\ensuremath{ \mathbf y}}
\newcommand{\mbz}{\ensuremath{ \mathbf z}}


\DeclareMathOperator{\card}{\text{card}}
\DeclareMathOperator{\conv}{\text{conv}}
\DeclareMathOperator{\diam}{\text{diam}}
\DeclareMathOperator{\height}{\text{ht}}
\DeclareMathOperator{\id}{id}
\DeclareMathOperator{\length}{length}
\DeclareMathOperator{\orbit}{\text{orbit}}
\newcommand{\bmat}{\left(\begin{matrix}}
\newcommand{\emat}{\end{matrix}\right)}
\newcommand{\bsmallmat}{\left(\begin{smallmatrix}}
\newcommand{\esmallmat}{\end{smallmatrix}\right)}
\DeclareMathOperator{\qand}{\quad\text{and}\quad}
\DeclareMathOperator{\qqand}{\qquad\text{and}\qquad}

\DeclareMathOperator{\support}{\text{support}}
\DeclareMathOperator{\image}{\text{image}}
\DeclareMathOperator{\Perm}{\text{Perm}}
\DeclareMathOperator{\Tet}{\text{Tet}}

\DeclareMathOperator{\vectoran}{\left( \begin{matrix} a_1 \\ \vdots \\ a_n \end{matrix}\right)}
\DeclareMathOperator{\vectorbn}{\left( \begin{matrix} b_1 \\ \vdots \\ b_n \end{matrix}\right)}
\DeclareMathOperator{\vectorcn}{\left( \begin{matrix} c_1 \\ \vdots \\ c_n \end{matrix}\right)}
\DeclareMathOperator{\vectorrn}{\left( \begin{matrix} r_1 \\ \vdots \\ r_n \end{matrix}\right)}
\DeclareMathOperator{\vectorxn}{\left( \begin{matrix} x_1 \\ \vdots \\ x_n \end{matrix}\right)}
\DeclareMathOperator{\vectoryn}{\left( \begin{matrix} y_1 \\ \vdots \\ y_n \end{matrix}\right)}
\DeclareMathOperator{\vectorzn}{\left( \begin{matrix} z_1 \\ \vdots \\ z_n \end{matrix}\right)}
\DeclareMathOperator{\vectorxy}{\left( \begin{matrix} x \\ y \end{matrix} \right)}
\DeclareMathOperator{\vectorxyz}{\left( \begin{matrix} x \\ y \\z \end{matrix} \right)}
\DeclareMathOperator{\vectoram}{\left( \begin{matrix} a_1 \\ \vdots \\ a_m \end{matrix}\right)}
\DeclareMathOperator{\vectorbm}{\left( \begin{matrix} b_1 \\ \vdots \\ b_m \end{matrix}\right)}
\DeclareMathOperator{\vectorcm}{\left( \begin{matrix} c_1 \\ \vdots \\ c_m \end{matrix}\right)}
\DeclareMathOperator{\vectorrm}{\left( \begin{matrix} r_1 \\ \vdots \\ r_m \end{matrix}\right)}
\DeclareMathOperator{\vectorxm}{\left( \begin{matrix} x_1 \\ \vdots \\ x_m \end{matrix}\right)}
\DeclareMathOperator{\vectorym}{\left( \begin{matrix} y_1 \\ \vdots \\ y_m \end{matrix}\right)}
\DeclareMathOperator{\vectorzm}{\left( \begin{matrix} z_1 \\ \vdots \\ z_m \end{matrix}\right)}

\DeclareMathOperator{\vectorch}{\left( \begin{matrix} c_1 \\ \vdots \\ c_h \end{matrix}\right)}
\DeclareMathOperator{\vectorsmallch}{\left( \begin{smallmatrix} c_1 \\ \vdots \\ c_h \end{smallmatrix}\right)}


\DeclareMathOperator{\vectorzero}{\left( \begin{matrix} 0 \\ \vdots \\ 0 \end{matrix} \right)}
\DeclareMathOperator{\vectorone}{\left( \begin{matrix} 1 \\ \vdots \\ 1 \end{matrix} \right)}

\DeclareMathOperator{\vectorsmallan}{\left( \begin{smallmatrix} a_1 \\ \vdots \\ a_n \end{smallmatrix}\right)}
\DeclareMathOperator{\vectorsmallbn}{\left( \begin{smallmatrix} b_1 \\ \vdots \\ b_n \end{smallmatrix}\right)}
\DeclareMathOperator{\vectorsmallcn}{\left( \begin{smallmatrix} c_1 \\ \vdots \\ c_n \end{smallmatrix}\right)}
\DeclareMathOperator{\vectorsmallrn}{\left( \begin{smallmatrix} r_1 \\ \vdots \\ r_n \end{smallmatrix}\right)}
\DeclareMathOperator{\vectorsmallun}{\left( \begin{smallmatrix} u_1 \\ \vdots \\ u_n \end{smallmatrix}\right)}
\DeclareMathOperator{\vectorsmallvn}{\left( \begin{smallmatrix} v_1 \\ \vdots \\ v_n \end{smallmatrix}\right)}
\DeclareMathOperator{\vectorsmallwn}{\left( \begin{smallmatrix} w_1 \\ \vdots \\ w_n \end{smallmatrix}\right)}
\DeclareMathOperator{\vectorsmallxn}{\left( \begin{smallmatrix} x_1 \\ \vdots \\ x_n \end{smallmatrix}\right)}
\DeclareMathOperator{\vectorsmallyn}{\left( \begin{smallmatrix} y_1 \\ \vdots \\ y_n \end{smallmatrix}\right)}
\DeclareMathOperator{\vectorsmallzn}{\left( \begin{smallmatrix} z_1 \\ \vdots \\ z_n \end{smallmatrix}\right)}
\DeclareMathOperator{\vectorsmallxy}{\left( \begin{smallmatrix} x \\ y \end{smallmatrix} \right)}
\DeclareMathOperator{\vectorsmallxyz}{\left( \begin{smallmatrix} x \\ y \\z \end{smallmatrix} \right)}

\DeclareMathOperator{\vectorsmallam}{\left( \begin{smallmatrix} a_1 \\ \vdots \\ a_m \end{smallmatrix}\right)}
\DeclareMathOperator{\vectorsmallbm}{\left( \begin{smallmatrix} b_1 \\ \vdots \\ b_m \end{smallmatrix}\right)}
\DeclareMathOperator{\vectorsmallcm}{\left( \begin{smallmatrix} c_1 \\ \vdots \\ c_m \end{smallmatrix}\right)}
\DeclareMathOperator{\vectorsmallrm}{\left( \begin{smallmatrix} r_1 \\ \vdots \\ r_m \end{smallmatrix}\right)}
\DeclareMathOperator{\vectorsmallxm}{\left( \begin{smallmatrix} x_1 \\ \vdots \\ x_m \end{smallmatrix}\right)}
\DeclareMathOperator{\vectorsmallym}{\left( \begin{smallmatrix} y_1 \\ \vdots \\ y_m \end{smallmatrix}\right)}

\DeclareMathOperator{\vectorsmallzero}{\left( \begin{smallmatrix} 0 \\ \vdots \\ 0 \end{smallmatrix} \right)}
\DeclareMathOperator{\vectorsmallone}{\left( \begin{smallmatrix} 1 \\ \vdots \\ 1 \end{smallmatrix} \right)}

\DeclareMathOperator{\vector02}{\left( \begin{smallmatrix} 0 \\ 0\end{smallmatrix} \right)}
\DeclareMathOperator{\vector03}{\left( \begin{smallmatrix} 0 \\ 0 \\ 0 \end{smallmatrix} \right)}



\title[Sumset sizes in additive number theory]
{Triangular and tetrahedral number differences of sumset sizes 
in additive number theory}
\author{Melvyn B.  Nathanson}
\address{Department of Mathematics\\Lehman College (CUNY)\\
Bronx, NY 10468}
\email{melvyn.nathanson@lehman.cuny.edu}


%\dedicatory{To Carl Pomerance on his 80th birthday}

\date{\today}

\subjclass[2000]{11B05, 11B13, 11B30,11B34, 11B75} 

\keywords{Sumset,  distribution of of sumset sizes, 
triangular numbers, tetrahedral numbers, 
popular sizes, Sidon set, $B_h$-set,   
additive number theory, combinatorial number theory}
\thanks{Supported in part by  PSC-CUNY Research Award Program grant 66197-00 54.}


\begin{document}

 
\begin{abstract}
The study of sums of finite sets of integers has mostly concentrated on sets 
with very small sumsets (Freiman's theorem and related work) 
and on sets with very large sumsets (Sidon sets and $B_h$-sets).  
This paper considers the full range of sumset sizes of finite sets 
of integers and an unexpected 
pattern (related to the triangular and tetrahedral numbers) 
that appears in the distribution of 
popular sumset sizes of sets of size 4. 
\end{abstract}


\maketitle



\section{The sumset size problem}

The \emph{integer interval} $[u,v]$ is the set of all integers $n$ 
such that $u \leq n \leq v$. 


The $h$-fold sum of a set $A$ of integers is the set $hA$ of all sums of $h$ 
not necessarily distinct elements of $A$: 
\[
hA = \left\{ a_1+\cdots + a_h: a_i \in A \text{ for all } i \in [1,h] \right\}.
\]
It is a fundamental unsolved problem in additive number theory, and a problem 
that strangely seems not to have been previously investigated, 
to compute and understand the sumset sizes of finite sets 
of integers.\footnote{This 
 problem was discussed at three recent number theory conferences 
  (Integers 2025 in Athens, Georgia, CANT 2025 in New York, and BIRS-IMAG  in Granada, Spain) and no one could cite earlier work.  
  References to previous papers would be appreciated.}
For all positive integers $h$ and $k$, we define 
the \emph{range of sumset sizes}  
\[
\mcr_{\Z}(h,k) = \left\{ |hA|: A \subseteq \Z \text{ and } |A| = k \right\}.
\] 
This set was introduced in Nathanson~\cite{nath25bb} 
and studied in~\cite{nath25aa,nath2505,nath2505b}.  


The idea of ``sumset'' is central in additive number theory 
and suggests many simple but \emph{a priori} difficult questions.  
For example, does there exist a set $A$ of size $|A| = 11$ 
whose $7$-fold sum has size $|7A| = 15551$?

The answer is ``yes,'' and the set 
\[
A_0 = \{ 0,897, 2056,2441,2988,3259,5294,6506,8013,9391,9872\}
\] 
has the desired sumset sum.  However, the answer was obtained by cheating: I started with the set $A_0$ and  computed its $7$-fold sum.  
For this set $A_0$, Maple quickly calculated the sumsets $hA_0$ 
and the  sumset sizes $|hA_0|$ 
for $h \in [1,7]$.  The second line of the following table 
gives the sequence of sumset sizes $|hA_0|$.  
The maximum size of an $h$-fold sumset of a set of size $k$ 
is $\binom{h+k-1}{h}$. 
The third line of the table gives the maximum size of an $h$-fold sumset 
of a set $A$ of size 11 for $h \in [1,7]$. 

\begin{center}
\begin{tabular}{c|r|r|r|r|r|r|r}
$h$ & 1 & 2 & 3 &4 & 5 & 6 & 7 \\ \hline
$|hA_0|$& 11 & 66 & 285 &  986 & 2878 & 7226 & 15551 \\ \hline 
$\max |hA|$ &  11 & 66 & 286 & 1001 & 3003 & 8008 & 19448 
\end{tabular}
\end{center}

 A $B_h$-set is a set $A$ of integers such that every integer 
 in the $h$-fold sumset $hA$ has a unique representation as a sum of 
 $h$ increasing elements of $A$.  
 If $A$ has size $k$, then $A$ is a $B_h$-set if and only if 
 $|hA| =  \binom{h+k-1}{h}$. 
 A $B_2$-set is also called a Sidon set.
 We see from the table above that $A_0$ is a Sidon set.  
Every $B_3$-set of size 11 satisfies $|3A| = 286$, but the special set $A_0$ 
has $|3A_0| = 285$, and so is not a $B_3$-set.  
Moreover, there must be exactly one integer $n_0$ that has two 
representations as a sum of three elements of $A_0$ and 
no integer that has more than two representations 
as a sum of three elements of $A_0$.  We have 
\[
n_0 = 18782 = 0 + 9391 + 9391 = 897 + 8013 + 9872 \in 3A_0.
\]


There are many related questions.  
The maximum size of a 7-fold sumset of a set of 10 integers is 
$\binom{16}{7} = 11440$.  It follows that if $A$ is a set of integers 
with $|7A| = 15551$, then $|A| \geq 11$.  
What is the cardinality of the  largest set  $A$ such that $|7A| = 15551$?

Sets $A$ and $B$ are \emph{affinely equivalent} if there are numbers 
$\lambda \neq 0$ and $\mu$ such that 
\[
B = \lambda\ast A + \mu = \{\lambda a+ \mu: a \in A\}.
\]
Affinely equivalent sets have the same sumset sums.  
Sets $A$ and $B$ are \emph{affinely inequivalent} 
if they are not affinely equivalent.  
How many affinely inequivalent sets $A$ satisfy $|A| = 11$ 
and $|7A| = 15551$?
How many affinely inequivalent finite sets $A$ satisfy $|7A| = 15551$?
The computation complexity inequality~\eqref{complexity} 
implies that all these questions can be answered in finite time.    

We have 
\[
\max \mcr_{\Z}(h,k) = \binom{h+k-1}{k-1} 
\]
and sets with this sumset size ($B_h$-sets and Sidon sets) have been extensively studied (\cite{nath-ubiq,nath-Canada,nath2502,obry04}).
Similarly, 
\[
\min \mcr_{\Z}(h,k) = h(k-1)+1
\]
and sets with this sumset size (called arithmetic progressions) 
and sets with very small sumset size have also been extensively studied 
(Freiman's theorem and related work~\cite{nath1996bb,tao-vu06}). 
Thus,
\beq      \label{R-interval}
\mcr_{\Z}(h,k) \subseteq \left[ h(k-1)+1, \binom{h+k-1}{k-1} \right].
\eeq

Because the set $\mcr_{\Z}(h,k)$ is finite, we can define the 
integer $N(h,k)$ as the smallest integer $N$ such that 
\[  
\mcr_{\Z}(h,k) = \{ |hA|: A \subseteq [0,N-1] \text{ and } |A| = k \}.
\]
The integer $N(h,k)$ determines the computational complexity of the 
sumset size problem.  It is proved in~\cite{nath2505b} 
that, for all $h \geq 3$ and $k \geq 3$,
\beq         \label{complexity}
N(h,k) < 4(8h)^{k-1}.
\eeq
Thus, for fixed $k$, the set $\mcr_{\Z}(h,k)$ can be computed in polynomial time. 


Here are some simple results. 
For all $h$ and $k$, 
\[
\mcr_{\Z}(h,1) = \{1\}, \quad \mcr_{\Z}(1,k) = \{k\}, 
\qand
\mcr_{\Z}(h,2) = \{h+1\}.
\]
For all $k$, the set $\mcr_{\Z}(2,k)$ is the interval of integers:
\[
\mcr_{\Z}(2,k) = \left[ 2k-1, \binom{k+1}{2} \right]. 
\]
However, the set $\mcr_{\Z}(3,3)$ is not an interval of integers. We have 
\[
\mcr_{\Z}(3,3) \subseteq \left[ 7, 10 \right] 
\]
but 
\[
\mcr_{\Z}(3,3) = \{7,9,10\}.
\]
The number 8 missing.  For $k=3$ and all $h \geq 1$, 
Nathanson~\cite{nath25bb} proved that  
\[
\mcr_{\Z}(h,3) = \left\{  \binom{h+2}{2}  -  \binom{t}{2} : t \in [1,h] \right\}. 
\]
and so $|\mcr_{\Z}(h,3)| = h$. 


\bprob
For  $k=4$, we have 
\[
\mcr_{\Z}(h,4) \subseteq \left[ 3h+1, \binom{h+3}{3} \right].
\]
Compute $\mcr_{\Z}(h,4)$ for all positive integers $h$.   
\eprob



\bprob
For all positive integers $h$ and $k$,  
compute the set of sumset sizes $\mcr_{\Z}(h,k)$.  
Explain the  integers ``missing'' from the interval 
on the right side of~\eqref{R-interval}.  
This is a core problem in additive number theory.  
\eprob




\bprob
 Let $k \geq 4$.  
Prove that ``most'' numbers in the interval 
on the right side of~\eqref{R-interval} 
are \emph{not} sumset sizes, that is, 
prove that 
\[
\left| \mcr_{\Z}(h,k) \right|= o_k \left( h^{k-1} \right)
\]
or, better, that 
\[
\left| \mcr_{\Z}(h,k) \right|= O_k\left( h^{k-2} \right). 
\]
\eprob

\bprob
Is there some kind of relation or reciprocity 
between the sets $\mcr_{\Z}(h,k)$ 
and $\mcr_{\Z}(k,h)$?  For example, is there a function $f: \N \rightarrow \N$ 
such that, if $t \in \mcr_{\Z}(h,k)$, then $f(t) \in \mcr_{\Z}(k,h)$? 
Is there a function $g: \N \rightarrow \N$ 
such that $t \in \mcr_{\Z}(h,k)$ if and only if $g(t) \in \mcr_{\Z}(k,h)$? 
Are there  functions $\alpha,\beta,\gamma$ such that, 
if $t \in \mcr_{\Z}(h,k)$, then $\gamma(t) \in \mcr_{\Z}(\alpha(h), \beta(k))$? 
\eprob

For related work on ``sumset size races,'' 
see~\cite{fox-krav-zhan25,krav25,peri-roto25}. 





\section{5-fold sums of sets of size 4}
In this section we describe some computational experiments on 
5-fold sumsets of sets of size 4.  The calculations were programmed in Maple. 

\textbf{Experiment 1.} 
Let $q = 100$. 
There are 
\[
\binom{100}{4} = 3,921,225
\]
subsets of size 4 contained in the integer interval $[0,q-1]$.  
Their 5-fold sumsets are contained in $[16,56]$.  
For all $t \in [16,56]$, let $S(t)$ denote the number of sets 
$A \subseteq [0,q-1] = [0,99]$ such that  $|A| = 4$ and $|5A| = t$.  
The following table records the 
distribution of  sumset sizes.

\begin{center}
\begin{tabular}{cr|cr|cr} 
t & $S(t)$  & t & $S(t)$   & t& $S(t)$   \\ \hline
16& 1617 & 30&  4655 & 44 & 10574   \\
 17 & 0 & 31 & 3444& 45 & 18862 \\
 18 & 0 & 32 & 5273 & 46 & 450320 \\
 19 & 0 & 33 & 1906 &  47 &6706 \\
 20 & 2400  &  34& 5846 & 48& 6120  \\                      
21 & 1200  & 35 & 5400  &   49 &  0 \\
22 &0  & 36  & 280850  & 50 & 15002\\
23 & 0 & 37 & 1152  &  51 & 27332 \\
24& 4750  & 38& 3176  & 52 &  460734  \\
25 &  0 & 39  & 4612  &  53 & 7628 \\
26 &950 & 40  &  0 &  54 & 30688 \\
27 &4704 & 41  & 10012  &  55 & 667836 \\
28 & 0 & 42  & 5888  &  56 & 1867410 \\
29&  1568 &  43 &  2610 & &      \\
\end{tabular}
\end{center}
The relation $S(16) = 1617$ is equivalent to the statement that the interval 
$[0,99]$ contains exactly 1617  arithmetic progressions of length 4. 
The relation $S(56) = 1867410$ means that the interval 
$[0,99]$ contains exactly 1867410 $B_5$-sets.  

\textbf{Experiment 2.}
Let $q = 1,000$. 
Maple randomly chose 150,000 sets of size 4 
in the interval $[0,q-1]$ and computed their 5-fold sumsets.  

 The set of  experimentally observed sumset sizes is 
\begin{align*}
\{20, 27, 30,  32, 34, 35,  36,    39, 41, 42, 
  44, 45, 46, 47,   50, 51, 52, 53, 54, 55, 56\}. 
\end{align*}
The following table lists these elements $t$ of the set $R(5,4)$ and the number $S(t)$ of sets $A$  in the sample space with $|5A|= t$: \\

\begin{center}
\begin{tabular}{cr|cr}
$t$ & $S(t)$  & $|5A|$ & $S(t)$   \\ \hline
20 & 1 & 45 & 6 \\  
27 & 2 & 46 & 1738   \\                      
30 & 2  & 47& 1  \\
32& 2 &   50 & 6 \\
34& 2 &     51& 7   \\
35 & 1 &    52& 1957   \\
36& 982 &   53 & 2     \\
39& 2  & 54& 9  \\
41& 2 &  55 & 3251 \\
42 & 1 & 56& 142020 \\
44& 6 &    &  
\end{tabular}
\end{center}

 
 \textbf{Experiment 3.}
Let $q = 10,000$. 
Maple randomly chose 1,000,000 sets of size 4 
in the interval  $[0,q-1]$  and computed their 5-fold sumsets.  
 The set of  experimentally observed sumset sizes is  
\begin{align*}
\{21,36, 38, 39, 42, 44, 45, 46, 50, 51, 52, 54,  55, 56\}. 
\end{align*}
Let $q = 10,000$. 
The following table lists these elements $t$ of the set $R(5,4)$ 
and the number $S(t)$ of sets $A$  in the sample space with $|5A|= t$: \\

\begin{center}
\begin{tabular}{cr|cr}
$t$ & $S(t)$  & $t$ & $S(t)$ \\ \hline
21 & 584 & 46& 1429  \\
36 & 805 & 50 & 1 \\
38 & 1 &  51 & 1 \\
39 & 1 & 52 & 1654 \\
42 & 1 & 54 & 1 \\
44 &  1&  55 & 2559 \\
45 & 1 &  56 & 992961  
\end{tabular}
\end{center}

\textbf{Experiment 4.}
Let $q = 1,000,000$. 
Maple randomly chose 250,000 sets of size 4 
in the interval  $[0,q-1]$  and computed their 5-fold sumsets.  
 The set of  experimentally observed sumset sizes is  
\begin{align*}
\{ 36,  46,  52,   55, 56\}. 
\end{align*}
The following table lists these elements $t$ of the set $R(5,4)$ and the number $S(t)$ of sets $A$  in the sample space with $|5A|= t$: \\

\begin{center}
\begin{tabular}{cr }
$t$ & $S(t)$    \\ \hline
36 & 1 \\
 46& 4 \\
52 & 2\\
 55 &4 \\
 56 & 249989  
\end{tabular}
\end{center}

\bprob
In all four experiments, the sets with sumset size 56 predominate 
and these are the $B_5$-sets.  
The most popular sumset sizes are 56, 55, 52, 46, and 36. 
Their differences are the consecutive triangular numbers 1,3, 6, and 10.  
Why? 
\eprob


\bprob
Let $N = N(5,4) = 4\cdot 40^3 = 256,000$. Inequality~\eqref{complexity} 
implies that $\mcr_{\Z}(5,4)$ is the set of all 
5-fold sums of  4-element sets of integers contained 
in the interval $[0,N-1]$. 
Compute $\mcr_{\Z}(5,4)$. 

\eprob




\section{A mystery about popular sumset sizes}
Maple was used to generate random sets of 4 integers contained in the intervals $[0,999]$ and $[0,9999]$, to calculate 
their $h$-fold sumsets for $h \in [2,9]$, and to find the ``popular'' sumset sizes.  We obtained the following results.    \\


\begin{center}
\begin{tabular}{c|r|r}
$h$  & most popular sumset sizes & successive differences  \\ \hline
2 &      $   9, 10  $ & 1 \\ 
3 &      $    16,19, 20 $ &  3,1\\            
4 &     $    25,  31, 34, 35  $ & 6,3,1 \\
5 & $ 36, 46, 52, 55, 56 $ & 10, 6,3,1  \\
6 &  $49, 64, 74, 80, 83, 84 $& 15, 10, 6,3,1 \\
7 &   $ 64, 85, 100, 110, 116, 119, 120 $ & 21, 15, 10, 6,3,1   \\
8 & $   81, 109, 130, 145, 155, 161, 164, 165  $   & 28, 21, 15, 10, 6,3,1     \\
9 &     $   100, 136, 164, 185, 200, 210, 216, 219, 220 $   & 36, 28, 21, 15, 10, 6,3,1 \\  
\end{tabular}
\end{center}

\bprob
For all $h \in [2,9]$, the sumset size set $\mcr(h,4)$
 has $h$ popular sizes and the differences between consecutive 
 popular sumset sizes are the first $h-1$ 
triangular numbers.  Why? 
\eprob

Recall that the \emph{$j$th tetrahedral number} 
$\Tet_j = \binom{j+2}{3}$ is the sum of the first $j$ triangular numbers.  
The popular sumset sizes appear to be 
the differences between  $\binom{h+3}{3}$, 
which is the size of a 4-element 
$B_h$-set, and the $h$ consecutive tetrahedral numbers 
$\Tet_0,  \Tet_1,\ldots,  \Tet_{h-1}$. 

\bprob
Let $h \geq 3$.  
Prove that 
\[
\left\{ \binom{h+3}{3} - \Tet_j : j \in [1,h-1]\right\}  
\subseteq \mcr_{\Z}(h,4). 
\]
\eprob

\bprob
Let $h$ and $q$ be positive integers.  
Let 
\[
\mca(4,q) = \left\{ A \subseteq [0,q-1]: |A| = 4  \right\} 
\]
\[
\mca^*(4,q,h) = \left\{ A \in \mca(4,q):
\text{$A$ is not a $B_h$-set} \right\} 
\]
and 
\[
\mcb^*(4,q,h ) = \left\{ A \in \mca^*(4,q,h ) :
 |hA| \in \left\{ \binom{h+3}{3} - \Tet_j : j \in [1,h-1]\right\}  \right\} 
\]
Almost all 4-element subsets of $[0,q-1]$ are $B_h$-sets (Nathanson~\cite{nath-ubiq,nath-Canada}):
\[
\lim_{q\rightarrow \infty} \frac{| \mca^*(4,q,h) |}{|\mca(4,q) |} = 0.
\]
Is it true that the sumsets sizes of almost all sets $A \in \mca(4,q)$ that 
are not $B_h$-sets  (that is, sets in $\mca^*(4,q)$) are the popular sumset sizes?  Equivalently, is it true that 
\[
\lim_{q\rightarrow \infty} \frac{| \mcb^*(4,q,h) |}{|\mca^*(4,q,h) |} = 1.
\]
\eprob




\bprob
Discover the pattern of popular sumset sizes for 
sets of size $k \geq 5$. 
\eprob




\section{A toy problem for physics} 
For positive integers $h$ and $k$, let $\mcx_{h,k}$ be the 
set of $k$-tuples  $\mbx = (x_1,\ldots, x_k)$ of nonnegative integers 
such that $\sum_{i=1}^k x_i = h$.  
Let $A = \{a_1,\ldots, a_k\}$ be a set of $k$ integers 
and let $\vec{A} = (a_1,\ldots, a_k)$ be the associated lattice point 
in $\Z^n$.  
For all $\mbx = (x_1,\ldots, x_k) \in \mcx_{h,k}$, we define 
\[
\mbx\cdot \vec{A} = \sum_{i=1}^k x_ia_i. 
\]
Then 
\beq                \label{hAX}
hA = \left\{\mbx\cdot \vec{A}: \mbx \in \mcx_{h,k} \right\}. 
\eeq
Because $\mcx_{h,k}$ is invariant with respect to the permutation 
group $S_n$ (that is, $\mbx = (x_1,\ldots, x_k) \in \mcx_{h,k}$  
if and only if $\sigma \mbx = (x_{\sigma(1)},\ldots, x_{\sigma(k)}) \in \mcx_{h,k}$ 
for all $\sigma \in S_n$),  relation~\eqref{hAX} 
is independent of the ordering of the coordinates in the vector $\vec{A}$. 



The \emph{representation function} $r_{A,h}(n)$ counts the number 
of representations of $n$ as a sum of $h$ elements of $A$, that is, 
the number of $k$-tuples $\mbx = (x_1,\ldots, x_k) \in \mcx_{h,k}$ 
such that $\mbx \cdot \vec{A} = \sum_{i=1}^k x_ia_i = n$.  
Analogous to collisions of molecules in a gas, 
$k$-tuples $\mbx  \in \mcx_{h,k}$ and $\mby  \in \mcx_{h,k}$ 
``collide'' with respect to the set $A$ if the $h$-fold sums 
$\mbx\cdot \vec{A} $ and $\mby \cdot \vec{A}$ are equal.  
Moreover, like multiple collisions of molecules, 
multiple collisions of vectors in $\mcx_{h,k}$, that is, integers $n \in hA$ 
such that $r_{A,h}(n) \geq 3$, should  usually 
be much rarer than simple collisions with $r_{A,h}(n) = 2$, 
but not always.   For example, with $k \geq 5$, 
the integer interval $A = [0,k-1]$ 
satisfies $2A = [0, 2k-2]$ and  
\[
r_{A,2}(n) = \begin{cases}
\lfloor \frac{n}{2}\rfloor + 1& \text{if $0 \leq n \leq k-1$}\\
\lfloor \frac{2k-2-n}{2}\rfloor + 1& \text{if $k-1 \leq n \leq 2k-2$.} 
\end{cases}
\]
Thus,
\begin{align*}
r_{A,2}(n) & = 1      \qquad  \text{if $n \in \{ 0,1,2k-3,2k-2 \}$} \\
r_{A,2}(n) & =  2     \qquad  \text{if   $n \in \{ 2,3 ,2k-5,2k-4 \}$} \\  
r_{A,2}(n) & \geq 3 \qquad   \text{if  $4 \leq n \leq 2k-6$  } 
\end{align*}
and so 
\beq       \label{r3>r2}
\card\left( n \in \Z: r_{A,h}(n) \geq 3 \right) 
> \card\left( n \in \Z: r_{A,h}(n) = 2 \right)
\eeq
for all $k \geq 7$. 
Inequality~\eqref{r3>r2} should be atypical.  
A quantitative version of this remark is the following.

\bprob
For every sufficiently large integer $q$, let 
\[
\mca(k,q )= \left\{ A \subseteq [0,q-1]: |A| = k\right\}.
\]
Prove that 
\[
\card\left( n \in \Z: r_{A,h}(n) \geq 3 \right) 
< \card\left( n \in \Z: r_{A,h}(n) = 2 \right) 
\]
for almost all $A \in \mca(k,q)$. 
\eprob



Computing the patterns of sumset sizes of finite sets of integers can be considered as a toy problem for certain much deeper questions in 
mathematical physics, such as the Hilbert program of deducing 
macroscopic fluid mechanical phenomena from 
mesoscopic  and microscopic motions.  
A significant part  of recent work of Yu Deng,  
Zaher Hani, and Xiao Ma~\cite{deng-hani-ma25,slom25} 
on Hilbert's sixth problem (the derivation of the Navier-Stokes equation 
from the Boltzman  equation from elastic collisions of a gas) 
is combinatorial and has a number theoretical flavor.  
It would be of interest to examine the combinatorial results of 
Deng, Hani, and Ma, and and the similarities between phenomena 
of combinatorial additive number theory and fluid dynamics. 



\emph{Acknowledgements.}
I thank Christian Elsholtz, Harald Helfgott, Noah Kravitz, Kevin O'Bryant, 
and Steven Senger for helpful remarks on this paper. 

\begin{thebibliography}{99}

\bibitem{deng-hani-ma25}
Yu Deng,  Zaher Hani, and Xiao Ma, 
Hilbert's sixth problem: derivation of fluid equations via Boltzmann's kinetic theory, 
arXiv.2503.01800.

\bibitem{fox-krav-zhan25}
J. Fox, N. Kravitz, and S. Zhang,
Finer control on relative sizes of iterated sumsets,
preprint, 2025. 


\bibitem{krav25}
N. Kravitz, Relative sizes of iterated sumsets, 
J. Number Theory 272 (2025), 113--128, 



\bibitem{nath1996bb}
M. B. Nathanson, 
\emph{Additive Number Theory: Inverse Problems and the Geometry of 
Sumsets}, Springer-Verlag, New York, 1996. 

\bibitem{nath-ubiq}
M.  B. Nathanson, On the ubiquity of Sidon sets, 
in: \emph{Number Theory (New York, 2003)},
Springer, New York, 2004, pages 263--272.

\bibitem{nath-Canada}
M.  B. Nathanson, $B_h$-sets of real and complex numbers, 
Canad. Math. Bull.  (2025),  to appear.
arXiv:2502.21272  


\bibitem{nath25bb}
M.  B. Nathanson, Problems in additive number theory, VI:  
Sizes of sumsets of finite sets, Acta Math. Hungar., to appear.
arXiv:2411.02365.



\bibitem{nath25aa}
M.  B. Nathanson, Inverse problems for sumset sizes of finite sets of integers, 
Fibonacci Quarterly, to appear. 
arXiv:2412.16154. 

\bibitem{nath2502}
M.  B. Nathanson,
${\mathbf Q}$-independence and the construction  of $B_h$-sets 
of integers and lattice points, 
arXiv:2502.13039.

\bibitem{nath2505}
M.  B. Nathanson, Explicit sumset sizes in additive number theory, 
arXiv:2505.05329. 


\bibitem{nath2505b}
M.  B. Nathanson, Compression and complexity for sumset sizes 
 in additive number theory, arXiv:2505.20998.

\bibitem{obry04}
K. O'Bryant, A complete annotated bibliography of work related to
  {S}idon sequences, Electron. J. Combin. DS11  (2004), 39.



\bibitem{peri-roto25}
P.  P\' eringuey and A. de Roton, 
A note on iterated sumsets races,
arXiv:2505.11233. 

\bibitem{schi25}
V. Schinina, On the sumset  of sets of size $k$, 
arXiv:2505.07679. 

\bibitem{slom25}
L. Sloman, Epic effort to ground physics in math opens up the secrets of time, 
Quanta Magazine, June 11, 2025.

\bibitem{tao-vu06}
T. Tao and V. Vu, \emph{Additive Combinatorics}, 
Cambridge Univ. Press, 2006.




\end{thebibliography}



\end{document}%%%%%%%%%%%%%%%%%%%%%%%%%%%%%%%%%%%%%%%%%%%%%%%%%%%%%%%%%%%%%%%%%%%%%%%%%%%%%%%%%%%%%%%%%%%%%%%%%%%%%%%%%%%%%%%%%%%%%%%%%%%%%%%%%%%%%%%%%%%%%%%%%%%%%%%%%%%%%%%%%%%%%%%%%%%%%%%%%%%%%%%%%%%%%%%%%%%%%%%%%%%%%%%%%%%%%%%%%%%%%%%%%%%%%%%%%%%%%%%%%%%%%%%%%%%%%%%%%%%%%%%%%%%%%%%%%%%%%%%%%%%%%%%%%%%%%%%%%%%%%%%%%%%%%%%%%%%%%%%%%%%%%%%%%%%%%%%%%%%%%%%%%%%%%%%%%%%%%%%%%%%%%%%%%%%%%%%%%%%%%%%%%%%%%%%%%%%%%%%%%%%%%%%%%%%%%%%%%%%%%%%%%%%%%%%%%%%%%%%%%%%%%%%%%%%%%%%%%%%%%%%%%%%%%%%%%%%%%%%%%%%%%%%%%%%%%%%%%%%%%%%%%%%%%%%%%%%%%%%%%%%%%%%%%%%%%%%%%%%%%%%%%%%%%%%%%%%%%%%%%%%%%%%%%%%%%%%%%%%%%%%%%%%%%%%%%%%%%%%%%%%%%%%%%%%%%%%%%%%%%%%%%%%%%%%%%%%%%%%%%%%%%%%%%%%%%%%%%%%%%%%%%%%%%%%%%%%%%%%%%%%%%%%%%%%%%%%%%%%%%%%%%%%%%%%%%%%%%%%%%%%%%%%%%%%%%%%%%%%%%%%%%%%%%%%%%%%%%%%%%%%%%%%%%%%%%%%%%%%%%%%%%%%%%%%%%%%%%%%%%%%%%%%%%%%%%%%%%%%%%%%%%%%%%%%%%%%%%%%%%%%%%%%%%%%%%%%%%%%%%%%%%%%%%%%%%%%%%%%%%%%%%%%%%%%%%%%%%%%%%%%%%%%%%%%%%%%%%%%%%%%%%%%%%%%%%%%%%%%%%%%%%%%%%%%%%%%%%%%%%%%%%%%%%%%%%%%%%%%%%%%%%%%%%%%%%%%%%%%%%%%%%%%%%%%%%%%%%%%%%%%%%%%%%%%%%%%%%%%%%%%%%%%%%%%%%%%%%%%%%%%%%%%%%%%%%%%%%%%%%%%%%%%%%%%%%%%%%%%%%%%%%%%%%%%%%%%%%%%%%%%%%%%%%%%%%%%%%%%%%%%%%%%%%%%%%%%%%%%%%%%%%%%%%%%%%%%%%%%%%%%%%%%%%%%%%%%%%%%%%%%%%%%%%%%%%%%%%%%%%%%%%%%%%%%%%%%%%%%%%%%%%%%%%%%%%%%%%%%%%%%%%%%%%%%%%%%%%%%%%%%%%%%%%%%%%%%%%%%%%%%%%%%%%%%%%%%%%%%%%%%%%%%%%%%%%%%%%%%%%%%%%%%%%%%%%%%%%%%%%%%%%%%%%%%%%%%%%%%%%%%%%%%%%%%%%%%%%%%%%%%%%%%%%%%%%%%%%%%%%%%%%%%%%%%%%%%%%%%%%%%%%%%%%%%%%%%%%%%%%%%%%%%%%%%%%%%%%%%%%%%%%%%%%%%%%%%%%%%%%%%%%%%%%%%%%%%%%%%%%%%%%%%%%%%%%%%%%%%%%%%%%%%%%%%%%%%%%%%%%%%%%%%%%%%%%%%%%%%%%%%%%%%%%%%%%%%%%%%%%%%%%%%%%%%%%%%%%%%%%%%%%%%%%%%%



 \textbf{Experiment x.}
 The set of  experimentally observed sumset sizes is 
\begin{align*}
\{24, 27, 30,  31, 32,   36,    39, 41, 42, 
  44, 45, 46, 47, 48, 50, 51, 52, 53, 54, 55, 56\}. 
\end{align*}
The following table lists these elements $t$ of the set $R(5,4)$ and the number $S(t)$ of sets $A$  in the sample space with $|5A|= t$: \\

\begin{center}
\begin{tabular}{cr|cr}
$t$ & $S(t)$  & $|5A|$ & $S(t)$   \\ \hline
24 & 2 & 46 & 2081 \\  
27 & 3 &  47& 2 \\                      
30 & 2  &  48& 5  \\
31& 2 & 50 & 5  \\
32& 3 & 51& 10    \\
36& 1273 &   52& 2323  \\
39& 2  &  53 & 3    \\
41& 2 &   54& 24 \\
42 & 4 & 55 & 3688 \\
44& 5 &  56& 140548  \\
45 & 6 & & \\
\end{tabular}
\end{center}



 \textbf{Experiment 3.}
Let $q = 10,000$. 
Maple randomly chose 1,000,000 sets of size 4 
in the interval  $[0,q-1]$  and computed their 5-fold sumsets.  
 The set of  experimentally observed sumset sizes is  
\begin{align*}
\{21,36, 38, 39, 42, 44, 45, 46, 50, 51, 52, 54,  55, 56\}. 
\end{align*}
The following table lists these elements $t$ of the set $R(5,4)$ and the number $S(t)$ of sets $A$  in the sample space with $|5A|= t$: \\

\begin{center}
\begin{tabular}{cr|cr}
$t$ & $S(t)$  & $t$ & $S(t)$ \\ \hline
21 & 584 & 46& 1429  \\
36 & 805 & 50 & 1 \\
38 & 1 &  51 & 1 \\
39 & 1 & 52 & 1654 \\
42 & 1 & 54 & 1 \\
44 &  1&  55 & 2559 \\
45 & 1 &  56 & 992961  
\end{tabular}
\end{center}









%%%%%%%%%%%%%%%%%%%%%%%%%%%%%%%%%%%%%%%%%%%%%%%%%%%%%%%%%%%%%%%%%%%%%%%%%%%%%%%%%%%%%%%%%%%%%%%%%%%%%%%%%%%%%%%%%%%%%%%%%%%%%%%%%%%%%%%%%%%%%%%%%%%%%%%%%%%%%%%%%%%%%%%%%%%%%%%%%%%%%%%%%%%%%%%%%%%%%%%%%%%%%%%%%%%%%%%%%%%%%%%%%%%%%%%%%%%%%%%%%



\begin{center}
\begin{tabular}{cc|cc|cc} 
$|5A|$ & number  & $|5A|$ & number & $|5A|$ & number  \\ \hline
16&  & 30&   & 44 &    \\
 17 & & 31 & & 45 \\
 18 & & 32 &  & 46 &  \\
 19 & & 33 &  &  47 & \\
 20 &   &  34&  & 48&   \\                      
21 &   & 35 &   &   49 &   \\
22 & & 36  &   & 50 & \\
23 & & 37 &   &  51 &  \\
24&   & 38&   & 52 &    \\
25 &   & 39  &   &  53 & \\
26 & & 40  &   &  54 & \\
27 & & 41  &   &  55 & \\
28 & & 42  &   &  56 & \\
29&   &  43 &  & &      \\
\end{tabular}
\end{center}


\vspace{2cm}

\begin{center}
\begin{tabular}{c|r}
$h$  & most popular sumset sizes   \\ \hline
2 &      $   9, 10  $  \\ 
3 &      $    16,19, 20 $  \\            
4 &     $    25,  31, 34, 35  $   \\
5 & $ 36, 46, 52, 55, 56 $   \\
6 &  $49, 64, 74, 80, 83, 84 $  \\
7 &   $ 64, 85, 100, 110, 116, 119, 120 $    \\
8 & $   81, 109, 130, 145, 155, 161, 164, 165  $       \\
9 &     $   100, 136, 164, 185, 200, 210, 216, 219, 220 $     \\  
\end{tabular}
\end{center}

\vspace{2cm}




%%%%%%%%%%%%%%%%%%%%%%%%%%%%%%%%%%%%%%%%%%%%%%%%%%%%%%%%%%%%%%%%%%%%%%%%%%%%%%%%%%%%%%%%%%%%%%%%%%%%%%%%%%%%%%%%%%%%%%%%%%%%%%%%%%%%%%%%%%%%%%%%%%%%%%%%%%%%%%%%%%%%%%%%%%%%%%%%%%%%%%%%%%%%%%%



Equivalently, the following table lists  elements $\ell$ of the set $R(5,4)$ and the fraction of sets $A$  in the sample space with $|5A|=\ell$: \\

                      16, 0.00002987928768 
                      
                     20, 0.000009959762559 
                     
                      21,  0.00001991952512
                      
                     24, 0.000009959762559

                     29, 0.000009959762559

                     30, 0.000009959762559

                     32, 0.000009959762559

                     34, 0.000009959762559

                     35, 0.000009959762559

                       36, 0.007967810047

                     38, 0.000009959762559

                      39, 0.00001991952512

                      41, 0.00001991952512

                     42, 0.000009959762559

                     43, 0.000009959762559

                      44, 0.00001991952512

                      45, 0.00006971833791

                       46, 0.01402334568

                      47, 0.00001991952512

                      50, 0.00004979881280

                      51, 0.00006971833791

                       52, 0.01590574081

                      53, 0.00002987928768

                      54, 0.00002987928768

                       55, 0.02491932592

                        56, 0.9367057089 \\


\section{$h=5$ and $k=4$}


\section{$h = 7$ and $k = 4$}


\noindent
22, 28, 29, 34, 36, 39, 41, 43, 44, 45, 47, 48, 49, 50, 52, 53, 54, 55, 56, 57, 58, 59, 60, 
62, 63, 64, 65, 66, 67, 68, 69, 70, 71, 72, 73, 74, 75, 77, 78, 79, 80, 81, 82, 83, 84, 85, 86, 87, 89, 90, 91, 92, 93, 94, 95, 96, 97, 98, 99, 100, 101, 102, 103, 104, 105, 106, 107, 108, 109, 110, 111, 112, 114, 115, 116, 117, 118, 119, 120


\section{Notes}
From email on May 31, 2025:

\bt[O'Bryant]
The number of 4-element Sidon sets in $[1,n]$ is 
\[
s_n = \frac{1}{24} n^4 - \frac{7}{12} n^3 + \frac 83 n^2+O(n)
\]
\et



Other:

Kohayakawa, Y., Lee, S.J., Rödl, V. and Samotij, W. (2015), 
The number of Sidon sets and the maximum size of Sidon sets contained 
in a sparse random set of integers. 
Random Struct. Alg., 46: 1-25. https://doi.org/10.1002/rsa.20496

\newpage

\section{Problems - June 1, 2025}

What is the fraction of $k$-element subsets of $[0,q-1]$ that are not $B_h$ sets? 

What is the fraction of $3$-element subsets of $[0,q-1]$ that are not $B_h$ sets? 

What is the fraction of $4$-element subsets of $[0,q-1]$ that are not $B_h$ sets? 

For $\ell \geq 3$, the number of $\ell$-length arithmetic progressions contained in 
$[0,q-1]$ is at most $\binom{q}{2}$. 

In the interval $[0,q-1]$, the number of solutions of 
\[
a+b = c+d
\]
is at most 



An arithmetic progression in $[0,q-1]$ is $\{a,a+d,a+2d\}$,
$a \in 0,q-1]$, $d \geq 1$, and $a+2d \leq q-1$.   
Thus, the number of 3-term arithmetic progressions in $[0,q-1]$ is 
\begin{align*}
\sum_{a=0}^{q-1} \left[ \frac{q-1-a}{2} \right] 
& = \sum_{a=0}^{q-3} \left[ \frac{q-1-a}{2} \right]  \\ 
& = \sum_{a=0}^{q-3} \frac{q-1-a}{2} + O(q) \\ 
& = \frac{1}{2} \sum_{a=2}^{q-1} a \\
& = \frac{q^2}{4} + O(q).
\end{align*}

Let $\lfloor x \rfloor$ denote the largest integer $n \leq x$ and 
$\lceil x \rceil$ the smallest integer $n \geq x$.  


\begin{lemma}
For $2x \in [2,2q-4]$,   
the number $N(x)$ of quadruples $(a,b,c,d)$  such that  
\beq              \label{abcd<}
0 \leq a < b \leq c < d \leq q-1
\eeq
and
\beq              \label{abcd}
a+d = b+c = 2x
\eeq 
is 
\[
N(x) = \begin{cases}
\binom{\lfloor x \rfloor+ 1}{2} & \text{if $x \leq (q-1)/2$}\\
\binom{q-\lceil x \rceil}{2} & \text{if $x \geq (q-1)/2$.}
\end{cases}
\] 
\end{lemma}


\begin{proof} 
If $x \leq (q-1)/2$, then for all $\{a,b\} \subseteq [0,\lfloor x \rfloor]$ 
with $a < b$ we have 
\[
0 \leq a < b \leq x \leq  2x-b = c < d = 2x-a \leq 2x \leq q-1
\]
and so the number of quadruples $(a,b,c,d)$  satisfying~\eqref{abcd<} 
and~\eqref{abcd} is $\binom{\lfloor x \rfloor + 1}{2}$.

If $x \geq (q-1)/2$, then  for all $\{c,d\} \subseteq [\lceil x \rceil, q-1]$ 
with $c < d$ we have 
\[
0 \leq q-1-d \leq 2x-d = a < b = 2x-c  \leq x \leq c < d  \leq q-1
\]
and so the number of quadruples satisfying~\eqref{abcd<} 
and~\eqref{abcd} $(a,b,c,d)$  is $\binom{q-\lceil x \rceil}{2}$.
This completes the proof. 
\end{proof}

\end{document}


\begin{problem}
Let $f_{h,k}(q)$ be the density of $k$-element subsets of $[1,q]$ that are $B_h$ sets.  
Is $f_{h,k}(q)$ increasing?
\end{problem}



The results are as follows:

\[
R(2,4) \subseteq [7,10]
\]

Frequent sumset sizes: 
\[
9, 10
\]


\[
R(  3,4) \subseteq [10, 20    ]
\]

Frequent sumset sizes: 
\[
16,19,20
\]



\[
R(  4,4) \subseteq [13, 35    ]
\]

Frequent sumset sizes: 
\[
25,  31, 34, 35
\]



\[
R(  5,4) \subseteq [ 16, 56   ]
\]

Frequent sumset sizes: 
\[
36, 46, 52, 55, 56
\]



\[
R( 6 ,4) \subseteq [ 19, 84   ]
\]

Frequent sumset sizes: 
\[
49, 64, 74, 80, 83, 84
\]



\[
R( 7 ,4) \subseteq [ 22, 120   ]
\]

Frequent sumset sizes: 
\[
64, 85, 100, 110, 116, 119, 120
\]



\[
R( 8 ,4) \subseteq [  25, 165  ]
\]

Frequent sumset sizes: 
\[
81, 109, 130, 145, 155, 161, 164, 165
\]



\[
R( 9 ,4) \subseteq [28,220    ]
\]

Frequent sumset sizes: 
\[
100, 136, 164, 185, 200, 210, 216, 219, 220
\]



\section{Special case $h=5$ and $k = 6$}
Choosing 526,852 random 4-element subsets of $[1,1000]$ generated the following data: 
\[
R_{\mathbf{Z}}(5,4) \subseteq [16,56]
\]
and 
\[
R_{\mathbf{Z}}(5,4) \supseteq [16,56] \setminus \{17,18,19,22,23,25,26,28,37,40,49\}
\]
The missing numbers 17, 18, 19 are covered by Vincent Schinina's theorem.

Problem:  Explain the missing numbers?

Problem:  Is it true that 
\[
R_{\mathbf{Z}}(5,4) = [16,56] \setminus \{17,18,19,22,23,25,26,28,37,40,49\} ?
\]

\newpage%%%%%%%%%%%%%%%%%%%%%%%%%%%%%%%%%%%%%%%%%%%%%%%%%%%%%%%%%%%%%%%%%%%%%%%%%%%%%%%%%%%%%%%%%%%%%%%%%%%%%%%%%%%%%%%%%%%%%%%%%%%%%%%%%%%%%%%%%%%%%%%%%%%%%%%%%%%%%%%%%%%%%%%%%%%%%%%%%%%%%%%%%%%%%%%%%%%%%%%%%%%%%%%%%%%%%%%%%%%%%%%%%%%%%%%%%%%%%%%%%%%%%%%%%%%%%%%%%%%%%%%%%%%%%



$R(7,3)$

\[
\{15, 21, 26, 30, 33, 35, 36\}
\]


$R(8,3)$
\[
\{17, 24, 30, 35, 39, 42, 44, 45\}
\]







$R(2,4) = \{7,8,9,10\}$

                      7, 0.00009668141058

                       8, 0.0001691924685

                        9, 0.02600729945

                        10, 0.9737268267

\[
\{9,10\}
\]

$R(3,4) = \{10, 12, 13, 14, 15, 16, 17, 18, 19, 20 \}$

\[
\{16,19,20\}
\]
                        10, 0.0003495316276

                      12, 0.0006091836939

                      13, 0.0002197055945

                       14, 0.001068568119

                       15, 0.001058581501

                       16, 0.07270257854

                       17, 0.001168434298

                       18, 0.001957377115

                        19, 0.1216170332

                        20, 0.7992490063
\\

$R(4,4) \supseteq \{23, 24, 25, 29, 30, 31, 34, 35\}$ 
\[
 \{25,  31, 34, 35\}
\]

                       25, 0.01406838667

                      29, 0.00004768944633

                      30, 0.0001907577853

                       31, 0.02432161763

                       34, 0.02856597835

                        35, 1.643998283


$R(5,4)$ 
\[
\{46, 52, 55, 56\}
\]


$R(6,4)$ 
\[
\{49, 64, 74, 80, 83, 84\}
\]
    49, 0.001273769448

                       64, 0.001455736512

                       74, 0.002365571831

                       80, 0.002456555363

                       83, 0.002456555363

                        84, 0.9899918115

$R(7, 4)$
               
               \[
               \{64, 85, 100, 110, 116, 119, 120 \}
\]

                       64, 0.0007621951220

                       85, 0.001286204268

                      100, 0.001524390244

                      110, 0.002524771341

                      116, 0.002096036585

                      119, 0.003334603659

                       120, 0.9884717988

 $R(8, 4)$
  \[
            \{81, 109, 130, 145, 155, 161, 164, 165 \}
\]
             
                       81, 0.001089324619

                      109, 0.001188354130

                      130, 0.002574767281

                      145, 0.002079619727

                      155, 0.001881560705

                      161, 0.004456327986

                      164, 0.003367003367

                       165, 0.9833630422


\section{Spectra for $h=9$ and $k=4$}

$R(9, 4)$
              \[
              \{100, 136, 164, 185, 200, 210, 216, 219, 220\}
              \]

                      

                       136, 0.01321095814

                       164, 0.01487971075

                       185, 0.02503128911

                       200, 0.02169378390

                       210, 0.03740787095

                       216, 0.03059379780

                       219, 0.04074537616

                       220, 0.8030871923


LL := seq([i, M94[i]], i = lower .. upper - 1);
plot([LL]);






Maple randomly chose N = 100,404 sets of size 4 
with elements in the interval of positive integers 
from 1 to 1000 and computed their 5-fold sumsets.  
 A set $A$ with $|A|=4$ and sumset size $|5A| = 16$ is a $4$-term arithmetic progression.   
   A set with sumset size 56 is a $B_5$-set.  
The set of sumset sizes experimentally observed is 
\begin{align*}
&\{16, 20, 21, 24, 29, 30,  32, 34,  36,  38, 39, 41, 42, \\
& \quad 43, 44, 45, 46, 47, 50, 51, 52, 53, 54, 55, 56\}
\end{align*}
%\{16, 20, 21, 24, 26, 27, 29, 30, 31, 32, 33, 34, 35, 36, 37, 38, 
% 39, 41, 42, 43, 44, 45, 46, 47, 48, 50, 51, 52, 53, 54, 55, 56\}
The following table lists  elements $\ell$ of the set $R(5,4)$ and the number of sets $A$  in the sample space with $|5A|=\ell$: \\

\begin{center}
\begin{tabular}{cc|cc|cc}% |cc|cc|cc}
$|5A|$ & number  & $|5A|$ & number & $|5A|$ & number  \\ % & $|5A|$ & number &$|5A|$ & number \\
\hline
16& 3 & 39& 2 \\  % & & & & &   \\
 20 & 1 &      41& 2 &  52& 1597 \\                      
21 & 2  & 42& 1 &      53 & 3 \\
24& 1 &    43& 1 & 54& 3  \\
29& 1 &  44& 2 &     55 & 2502 \\
30& 1 &  45& 7 &   56&  94049   \\
32& 1  &    46& 1408 \\
34& 1 &   47& 2 \\
36 &800 &  50& 5  \\
 38& 1 &     51& 7  \\
\end{tabular}
\end{center}

\bprob
In all three experiments, the sets with sumset size 56 predominate and these are the $B_5$-sets.  
The most popular sumset sizes are 56, 55, 52, 46, and 36. 
Their differences are the consecutive triangular numbers 1,3, 6, and 10.  
Why? 
\eprob



